\chapter[What Makes You the Same Person?]{What Makes Someone Still the Same Person?}
\section{The Photograph in the Drawer}
Start by picking up an object: an old photograph.

It probably lies at the bottom of a drawer at home, underneath a pile of certificates and instruction manuals. The child in it stands outside an elementary-school gate, wearing a uniform one size too large, missing a front tooth, smiling without a trace of self-consciousness. On the back, your mother has written: "First day of school."

That is you. Everyone says so, and you agree. But look at the photograph a little longer, and something strange begins to happen.

This child is much shorter than you, with softer hair; a permanent tooth now occupies that gap. You have had your eyesight checked at a hospital: you need glasses now, whereas the child did not. More importantly, a doctor will tell you that your body is constantly renewing itself: skin cells are replaced every few weeks, red blood cells roughly every four months. After several years, the vast majority of your cells are no longer the ones you had then. In terms of "materials" alone, you and the child in the photograph share very little.

Now look "inside." That child was most afraid of thunder, loved eating rice soaked in soup, and firmly believed in Santa Claus. Your fears, loves and beliefs have since changed several times over. Most things that child remembered, you no longer remember. Can you recall how you felt on the morning you started first grade? Probably not. Your connection through memory is as thin as a thread about to snap.

Yet nobody would seriously say, "That child no longer exists; this is a stranger." Your household registration has not changed; your name grew up with the child. Adults still hold you responsible for the neighbor's window that child broke. The law, your family and you yourself all insist: you are the same person.

On what grounds?

This question sounds pointless, like something philosophers invent when they have nothing better to do. Keep it in mind, though, because you will soon see that it is anything but pointless. It keeps appearing in hospital rooms, courtrooms, wills and phone accounts, forcing people to make real decisions they cannot avoid.

This chapter asks: \textbf{what makes someone still the same person?} Is a ship with every part replaced still the original ship? Is an older woman who has forgotten her entire past still herself? Is the "you" copied into a computer really you?

\section[A Granddaughter Called a Sister]{In the Hospital, She Calls Her Granddaughter Her Sister}
Now enter a scene.

It is three in the afternoon, in a neurology ward on the sixth floor of the inpatient building. The corridor smells of disinfectant mixed with rice porridge from the cafeteria. Sunlight cuts diagonally through the window at the far end, dividing the floor into light and shade. A call bell sounds intermittently at the nurses' station. An attendant pushes a treatment cart past the door; one wheel is slightly out of round and gives a little bump with every turn.

Grandma Chen, seventy-eight, sits on bed three. Her hair is neatly combed: her daughter Zhou Min did it that morning. She wears blue-and-white-striped hospital pajamas and holds an orange she has been squeezing for twenty minutes without peeling it.

Zhou Min sits in the visitor's chair beside the bed. Next to her stands her fourteen-year-old daughter, Xiaoman, still wearing her school uniform and backpack. Xiaoman comes to see her grandmother after school every Wednesday, a habit she has kept for almost two years.

"Mom, look who's here," Zhou Min says.

Grandma Chen looks up at Xiaoman, and her eyes brighten. She smiles and reaches out: "Little sister, you're here."

Little Sister was Grandma Chen's younger sister, who died more than thirty years ago.

Xiaoman's hand pauses on the edge of the bed before reaching forward for her grandmother to hold. She is very practiced at this movement now. She does not correct her grandmother, but says, "Yes, I'm here. Shall I peel your orange?"

Grandma Chen hands it over, then suddenly becomes suspicious again and lowers her voice: "That person in white just now---are they going to ask me for money? I don't have any. Your brother-in-law has all the money."

Zhou Min turns away to look out of the window.

This is not the first time this scene has happened, and it will not be the last. Grandma Chen has Alzheimer's disease, a disease that progressively damages the brain, often affecting recent memories first. She remembers her sister from sixty years ago but cannot recognize the granddaughter who visits every week. She remembers her husband managing their money when they were young but cannot recognize the daughter who has stayed beside her bed for two years. The doctor says the disease will continue to progress.

Last month, Xiaoman's uncle came back from his job in Shenzhen. The family talked in the small sitting room outside the ward, and it went badly. His exact words were: "Sis, this sounds harsh, but Mom isn't Mom anymore. The decisions she made, the people she cared about---they're gone. What's lying here now is just Mom's body. When we make decisions, we have to face reality."

Zhou Min immediately stood up. "She holds that orange waiting for my daughter to peel it. She remembers to leave the door open for Little Sister. How is she not Mom? You come home once a year. Of course she seems like a stranger to you."

On the surface, they were arguing over care arrangements and whether to sell the old house. But what made his words hurt lay underneath: \textbf{was their mother still their mother?} This was not a biological question: the body in the hospital bed was unquestionably Grandma Chen's. Nor was it a question of affection: both siblings loved their mother. It was a philosophical question entering a hospital room in everyday clothes: what must a person lose to cease being herself? What must remain for her still to be herself?

Remember this hospital room. Everything that follows tries to answer this question for that family---and for you.

\section{What Exactly Has Remained?}
Let us lay out the conflict before rushing to answer it.

We use "the same person" so effortlessly that we never inspect what it contains. Open it up now, and there are five candidate "batons": five things everyone considers important. The question is which one keeps "you" being "you."

\textbf{First baton: the body.} You remain you because your body has grown continuously from the child in the photograph to the present, without interruption. This baton sounds the sturdiest, but it has two cracks. First, as we have seen, most of the body's materials are replaced over several years: "the same body" plainly does not depend on the same atoms. Second, imagine an accident: someone loses both legs through injury or illness, has a heart valve replaced and receives artificial joints. Nobody says that this makes him a different person. Much of the body can be replaced while the judgment "still him" stays firmly in place.

\textbf{Second baton: memory.} You remain you because you remember yesterday and childhood; memory stitches together each day's version of you. But memory is precisely what failed first for Grandma Chen. If memory is the only baton, the uncle is right: the moment everything is forgotten, the person is gone. Most people cannot swallow that harsh conclusion. There is also a smaller problem: you remember nothing of yourself at eighteen months old. Was that a different person?

\textbf{Third baton: personality.} You remain you because your temperament, preferences and way of speaking are recognizable. But personality changes too. Serious illness can make a gentle person irritable; a long journey can make a timid person decisive. We all hear people say, "He's like a different person." Notice what that means: the person is still the same, but the personality has changed. A change of personality does not necessarily cancel out "the same person."

\textbf{Fourth baton: relationships.} You remain you because people know you, remember you and share a past with you. Grandma Chen remains a mother because the people she raised still stand beside her bed. This baton is comforting, but it also faces counterexamples. Someone stranded abroad and cut off from everyone, or an orphan whose entire family died in a disaster, has not thereby "become another person." A person living entirely alone remains himself.

\textbf{Fifth baton: a story.} You remain you because your life forms a line with a beginning and an end: birth, school, setbacks and changes of direction, its episodes echoing one another like a long novel. But stories are told, and can be told very differently. You can describe your life as "how I got where I am today," while someone else says, "He stopped being his old self long ago." The uncle and Zhou Min are really arguing over two different tellings.

Every one of these five batons can be dropped. Here is the real problem: \textbf{if everything can change, disappear or be replaced, what sustains "the same person"?} Or is "the same person" not sustained by any single thing at all?

Before reading on, take a position. Do you think Grandma Chen in that hospital room is still herself? Which baton supports your answer?

\section[The Full Case for Both Sides]{Giving Both Sides Their Full Case}
The argument in the sitting room outside the ward involved two positions. Here we will make the strongest complete case for each, including the one you disagree with. This book practices that move repeatedly: finish stating the other side's case before deciding whether to oppose it.

\textbf{Position one: when memory and mind break down, the person ends---the uncle's side.}

Let us complete his reasoning. It would be easy simply to condemn his words as "unfilial," but that would let the argument off too lightly.

First: what we love is the person, not the body. Someone is "Mom" because she remembers carrying you to the hospital when you had a fever, remembers which of her dishes you loved, remembers how she cried on your wedding day. Those memories and attachments are the entire content of "Mom." Now that content is disappearing: she calls her granddaughter her sister and treats her daughter as a stranger. The object of love is gone; what remains is a body needing care. Acknowledging this is not cruelty, but respect akin to respect for the dead. The real mother deserves mourning, not replacement by a shell impersonating her.

Second: refusing to acknowledge this leads to bad decisions. In middle-to-late-stage Alzheimer's, a patient may refuse treatment, fail to recognize herself in a mirror or demand to "go home" in the middle of the night. If we insist, "She is still herself; her wishes come first," then her current wishes---hiding an orange under her pillow, refusing a bath---must all be obeyed. The uncle thinks that only by acknowledging that "the decision-maker is no longer the original person" can relatives and doctors legitimately decide for her according to her \textbf{former} wishes. We will return to this reasoning in the section on medical decisions. It is not absurd.

Third: there is also an honest accounting of the cost. If everyone pretends that nothing has changed and stages "Grandma recognizes me" every week, who gets most exhausted? Children like Xiaoman. Saying the truth plainly hurts, but releases everyone from performing.

None of these three reasons is weak. Before you object, hold them in your hands and feel their weight.

\textbf{Position two: a person is much more than their memories---Zhou Min's side.}

Zhou Min's reasoning is equally complete. First: memory has never been the whole person. You remember nothing of yourself at eighteen months, but who would say you "weren't you yet"? An ordinary person forgets an elementary-school desk mate's name or what they ate a year ago today; nobody thinks part of them has died. Memory is some of a person's property, not the deed to the house. Grandma Chen has forgotten much, but she still worries whether there is enough money and saves an orange for "Little Sister." Her concern and thoughtfulness---things deeper than explicit memories---are still working.

Second: people live in relationships, not on hard drives. Half the meaning of "mother" resides in her attitude toward her children; the other half resides in their attitude toward her. Xiaoman comes every Wednesday, and Grandma Chen's eyes brighten: that instant shows the relationship remains. Her son comes home once a year, so naturally he sees a "stranger." But relationships are made through daily interaction, not remotely certified by blood ties. In Zhou Min's view, when he says, "Mom isn't Mom anymore," he reveals that \textbf{he himself} withdrew from the relationship first.

Third: declaring someone "already gone" is a dangerous move. Historically, defining a group as "no longer fully human" has almost always marked the beginning of something bad: slavery did it, and so did colonizers. These words in the hospital come from exhaustion, not malice, but their direction deserves scrutiny. Once we accept that "people without memories don't count as people," the next step is allowing them to be treated as less than people. Dignity should not be a prize awarded for passing a memory test.

Both sides have now had their say. Notice something: they do not dispute the facts. Both acknowledge that Grandma Chen cannot recognize her granddaughter; the evidence is not contested. They disagree about \textbf{what those facts mean}, a matter of interpretation and evaluation. Separating "what happened" from "what it means" is the first step in handling any disagreement. We will use that distinction repeatedly here.

Notice, too, that they are not competing over the same baton. The uncle stakes his case on memory and mind; Zhou Min on relationships and the personality still working underneath. A third baton has yet to enter: the body. In the next section, we will use a tool to test all three, one at a time.

\section[Thinking Tool: Thought Experiments]{Thinking Bridge: How to Run a Thought Experiment}
To examine questions such as "the same person," philosophers use \textbf{thought experiments}: constructing an arrangement in the mind that cannot be built in reality, and seeing what happens. The method has just three steps, all within your reach:

First, \textbf{change only one variable at a time}. When troubleshooting a machine, you cannot turn three screws at once, or you will not know which made the difference. Second, \textbf{watch your intuitive reaction}. After changing a variable, ask: is this still the original thing? Your intuition exposes the standard you actually believe in. Third, \textbf{find which theory gives way first}. Where intuitions clash, theories show their cracks.

Let us practice with three classic arrangements.

\textbf{Arrangement one: the Ship of Theseus.}

In the Life of Theseus in his Parallel Lives, the ancient Greek writer Plutarch recorded a dispute, paraphrased here. The Athenians preserved the hero Theseus's ship as a memorial, replacing each plank as it rotted. After many years, not one original plank remained. Philosophers then argued: was it still Theseus's ship?

Someone later made matters worse: suppose the discarded old planks were collected and reassembled into another ship. There would now be two ships: one entirely new in material but continuous in location, shape and use; the other entirely original in material but reassembled. \textbf{Which would be the Ship of Theseus?}

Set aside the answer and apply the three steps. There is just one variable: the speed of material replacement. If one plank is replaced each year, nobody doubts it is the same ship; if everything is replaced at once, intuition hesitates. What does that show? Your intuition is not using "identical materials" as its standard, but "continuous, gradual change, with the old structure carrying the new at every step." The bodily view scores a point: our bodies replace cells in just this way, changing most over several years, yet slowly and continuously enough for you to remain you.

\textbf{Arrangement two: the teletransporter.}

Now change another variable. Imagine a future teletransporter: you enter a booth in Beijing, a machine scans the position of every atom in your body and sends that information to Shanghai. The Shanghai booth assembles you exactly from new atoms, and the Beijing original is immediately destroyed. You experience only a moment of darkness before opening your eyes in Shanghai, with all your memories, your personality and the mole on your hand.

Would you use this teletransporter?

The British philosopher Derek Parfit seriously examined this arrangement in Reasons and Persons (1984). Intuitions split down the middle. Some say: of course; the person emerging remembers everything about me and thinks my thoughts, so he is me. Others say: never; the "me" in Beijing has been killed, and the person in Shanghai is merely an extraordinarily close copy. He \textbf{thinks} he is me, but that and actually being me are different things.

Notice what this clash reveals. The "continuity" standard that performed so well in the previous test has split into two: \textbf{physical continuity}, the gradual continuation of the same body, and \textbf{informational continuity}, the intact transmission of memories and personality. The teletransporter preserves the latter and destroys the former. Your decision whether to use it is a vote between those two kinds of continuity.

\textbf{Arrangement three: the divided brain.}

The hardest test comes next. Change the variable again: not copying, but dividing in two. Imagine that the two halves of your brain are transplanted into two new bodies. Ignore medical feasibility: a thought experiment need not be real, only clear. Both people wake up remembering your entire past, and both firmly believe they are you.

Where did the original person go? Say, "He became the person on the left"---why? What is lacking on the right? Say, "He became both"---but "identity" means a one-to-one relation; one person cannot simultaneously be two. Say, "Neither is him anymore; he died"---yet both lives have continued quite smoothly. What sort of death is that? Could there be a more comfortable death?

Parfit's conclusion from this test is famous. We will leave it hanging until "Deeper Challenge." For now, see the method's value: the three arrangements did not produce a final answer, but exposed the hidden standards in your intuitions. You thought you believed in "identity," but actually believed in "some kind of continuity." You thought there was only one continuity, but there are at least two, and they can conflict.

A thought experiment remains only a tool. It cannot vote for you, but it can help you know what you are voting for. Now let us look at an actual text from two thousand years ago and see how people approached the same question without experiments.

\section[Butterfly and Chariot]{Butterfly and Chariot: Two Ancient Testimonies}
Now read a short piece of primary material. More than two thousand years ago, people in China left a famous firsthand account concerning "whether I am still me."

The end of the Zhuangzi's "Discussion on Making All Things Equal" reads:

\begin{quote}
Once Zhuang Zhou dreamed he was a butterfly, a fluttering butterfly, delighted and perfectly at ease! He did not know he was Zhou. Suddenly he awoke, and there he was, unmistakably Zhou. He did not know whether Zhou had dreamed of being a butterfly, or a butterfly was dreaming of being Zhou. Between Zhou and the butterfly there must surely be a distinction. This is called the transformation of things.
\end{quote}
Roughly: Zhuang Zhou once dreamed he became a butterfly, fluttering happily, entirely unaware of being Zhuang Zhou. Suddenly he awoke, clearly Zhuang Zhou again. Then he was puzzled: had Zhuang Zhou dreamed of becoming a butterfly, or was a butterfly now dreaming of becoming Zhuang Zhou? There must be a distinction between the two, but what draws that boundary? Zhuangzi calls this "the transformation of things": all things transforming into one another.

Notice how Zhuangzi's question differs from ours. We asked, "How does a person remain himself?" He asks, "What makes us certain that this present self is the original?" In his dream, the "butterfly self" is equally complete and coherent, without any flaw. These passages sound playful, but deliver a serious blow. All our grounds for affirming "I am me" consist of present memories and feelings; the dreamer has a full set of those too. Each side possesses equally complete evidence.

Another source comes from ancient India. A Buddhist text called the Sutra of the Monk Nagasena, known in the Pali tradition as the Questions of King Milinda, records a conversation roughly two thousand years ago between Milinda, a king of Greek descent, and the Buddhist monk Nagasena. Its opening is wonderful; here is a paraphrase. The king asks Nagasena's name. Nagasena says "Nagasena" is only a name or designation: no substantial person called "Nagasena" exists. The king objects: if there is no such person, who is speaking and eating? Nagasena asks in return: you arrived by chariot, but what is the "chariot"? The wheels? The axle? The carriage? The shafts? None of them. If all the parts were dismantled and piled on the ground, where would the chariot be? The king admits that no individual part is the chariot; "chariot" is simply the designation people use when the parts are assembled. Nagasena replies: the same applies to a person. "Nagasena" is neither hair nor skin, neither memory nor temperament, but a convenient name people attach when these things work together.

This is a classic demonstration of Buddhist "not-self," called anattā in an ancient Indian language. Do not immediately object, "Isn't that just sophistry?" Its structure is exactly like your Ship of Theseus experiment: after every plank is replaced, where is "that ship"? After every part is removed, where is "that chariot"? Nagasena and ancient Greek philosophers encountered the same problem thousands of kilometers apart. That is not coincidence: the problem lies in the foundations of human reason, where anyone digging deeply enough will encounter it.

One warning, though: there is an especially easy mistake here, which we will unpack in "Deeper Challenge." Many readers instantly translate "not-self" into "you do not exist; everything about you is false." That translation is wrong, and harmfully so. We will examine why in that section.

\section[One Grandmother, Three Readings]{The Same Grandmother, Three Readings}
Back to the hospital room. We now have the pieces to construct three genuinely different readings of "whether Grandma Chen is still Grandma Chen." As usual, distinguish four things. What happened: she called her granddaughter her sister and held an orange waiting for someone to peel it; this is undisputed evidence. How the people there understood it: Zhou Min and her brother each had a view we have heard. What follows compares the second level, \textbf{why the judgment "she is still herself" can both hold and fail to hold}. How we should evaluate this and what we should do belong to the final level, left for you at the chapter's end.

\textbf{Reading one: bodily continuity is enough.}

This reading says people are animals; "the same person" and "the same tree" are the same kind of question. A tree that loses half its branches to lightning remains that tree because its life processes continue in the same body without interruption. Grandma Chen's body has lived continuously from eighteen to forty to seventy-eight, without a break: she is Grandma Chen, full stop. The advantage is clarity and verifiability: bodily continuity is visible and tangible, requiring no guessing. It also explains why we never doubt that the child in the photograph is you: that little body grew continuously into yours. Its limitation appeared in the teletransporter test. If identity depends only on bodily continuity, the person copied with all your information can never be you, even if he cares for your mother and remembers all your promises. Many cannot accept that. Besides, bodily continuity offers Zhou Min little help in the ward: it establishes that "Grandma Chen's body" is in the bed, without answering her brother's real question---whether \textbf{the person we love} is still there.

\textbf{Reading two: psychological continuity makes the person.}

This reading stakes its case on a chain connecting memories, personality and intentions. While the chain holds, the person remains; once it breaks sufficiently, the person comes apart. Its advantage is that it fits our feelings: we love someone's memories and temperament, not their liver cells. It also handles practical affairs particularly well. As we will see, social institutions mostly operate through psychological continuity in legal responsibility and medical authorization. It has two limitations. First is the question of degree: forgetting a school desk mate is unimportant, forgetting every relative matters, but where is the line between them? Who draws it? Alzheimer's worsens gradually; no day can be marked with a flag saying, "Since yesterday she has stopped being herself." Yet institutions and decisions often demand a definite date. The second limitation is sharper: the divided-brain test. Psychological continuity can become one-to-two, but identity cannot. Sooner or later, they part company.

\textbf{Reading three: people are woven from relationships and stories.}

This reading changes the viewpoint. "The same person" is not an object in a brain or body, but a fabric woven from other people's memories, names, promises and shared days. Grandma Chen remains a mother because the people she raised still treat her as a daughter treats a mother, and Xiaoman still appears every Wednesday as a granddaughter. Those actions themselves help constitute the statement "she is still herself." This reading explains why Zhou Min's and her brother's feelings are both real: someone weaving the fabric daily and someone inspecting it once a year naturally touch different fabrics. It also best accommodates the intuition that "people with memory loss still deserve human love." Its limitation is a dangerous dependence on the other weavers. What if everyone who remembers Grandma Chen is gone? Does she evaporate with those relationships? As we said, a person living alone remains himself. The relational view needs another piece of material---perhaps counting "the story a person tells herself" as a thread---to mend this hole.

Each reading holds part of the truth, and none reaches all of it. This is not a bland compromise. They actually capture three uses of "person": a person as an \textbf{animal}, involving the body; as a \textbf{subject}, involving psychology; and as a \textbf{role}, involving relationships and stories. Normally these uses overlap, so we never notice them. Extreme cases such as Alzheimer's, teletransporters and divided brains simply pull the three overlapping layers apart, allowing us to see for the first time that there are three.

Remember this three-layer structure. It is more durable than any answer consisting of a single slogan.

\section[Deeper Challenge: Beyond Identity]{Deeper Challenge: Perhaps Identity Is Not What Matters Most}
Now for Parfit's conclusion after the divided-brain test, one of the twentieth century's most powerful blows on this question.

Recall the arrangement: your brain is split in two and transplanted into two bodies. Two "yous" wake, each with all your memories and personality. Each of the three possible answers---you became the left-hand person, you became both, you died---has something wrong with it. Parfit says we cannot escape this dead end because we have asked the wrong question. We keep asking "whether the original person is still the same person," as though identity were a gold medal we must preserve. But look at what actually exists here: two complete lives continuing onward, with memories, personality, unfinished plans and love for the same people all continuing. None of these things is lost. The only "loss" is the abstract label "a one-to-one identity relation."

Parfit thus proposes a claim still debated throughout philosophy, roughly: \textbf{what matters is not identity, but psychological continuity and connectedness.} In other words, the "I am still me" we struggle to preserve may be only a shell. What is valuable inside it---someone existing tomorrow who remembers and continues today---still exists after division, even in duplicate. Parfit himself said that understanding this reduced his fear of death: if the hard boundary of "me" was never secure, then "my complete disappearance" no longer seemed so solidly terrifying.

Now bring in this chapter's second theory to sit opposite Parfit: Buddhist not-self. Earlier we flagged an easy misunderstanding; now we can unpack it. "The self does not exist" can mean at least three very different things. Confusing them causes serious trouble:

The first meaning: \textbf{nothing continuous exists, and all talk about "me" is empty.} This is annihilationism. Nagasena would not accept it. Although dismantling a chariot reveals no "chariot substance," it still carries goods, and responsibility still follows if it hits someone. Buddhist traditions likewise discuss causation and responsibility: what you do today has consequences for the future person who continues from you. If "not-self" meant "nothing exists and nobody is responsible for anything," Nagasena's conversation with the king would have been pointless.

The second meaning: \textbf{an independent, unchanging, all-controlling "self-substance" supposedly exists, but no such substance can be found.} This is what not-self actually seeks to say. It does not mean "you do not exist." You imagine an unchanging boss living inside your body, owning your hair, temperament and memories, remaining the same through life and death. But search the whole body, as you search every chariot part, and no boss appears. There are only changing, interacting processes: bodily metabolism, rising and falling sensations, memories forming and fading, habits strengthening and loosening. Buddhist traditions observe these in five categories---form, feeling, perception, formations and consciousness, collectively the "five aggregates." Roughly, "I" is where these five rivers meet, not an unmoving stone on the riverbed.

The third meaning: \textbf{"I" exists, but not in the way we think}. It is not a thing but an abbreviation, just as "chariot" abbreviates an arrangement of parts and "company" abbreviates people, contracts and an office. Nobody searches a warehouse for a substance called "company," yet companies can sign contracts, owe debts and be sued. Parfit's approach is not far from this: he accepts that people really exist, but argues that the answer to "what is a person?" contains no indivisible hard core.

Now consider their difference. Buddhist traditions take this observation toward spiritual practice: if "I" is a river rather than a stone, clinging to "me" is one source of suffering, and practicing release can make life less painful. Parfit moves toward ethics: if the boundaries of "me" are less rigid, the grounds for "caring only about myself" weaken. There is no wall labeled "identity" between your future self and my present self; similarly, perhaps the wall between me and others is thinner than I thought. One theory set out from India over two thousand years ago, the other from Oxford a few decades ago, and halfway along they exchanged a distant handshake.

Consider their limits too. The "river view" and "abbreviation theory" describe what people are well, but do less to tell you directly \textbf{what to do}. Should Grandma Chen's present wishes be respected in the ward? Should a court try an old person utterly different from the person of fifty years earlier? Theory illuminates the problem; people must still decide. A map is not a verdict.

\section[Yourself in Court, Ward and Phone]{The "You" in Court, in Hospital and on Your Phone}
This ancient question comes up for a hearing every day in three thoroughly modern places.

\textbf{The courtroom.} The law has the least room for ambiguity: punishment must fall on "the same person," not the wrong one. Its default answer is essentially bodily continuity plus a minimum of psychological continuity. A seventy-eight-year-old must answer for crimes committed at eighteen, even if personality, beliefs and memories have all changed. News of elderly war-crime defendants appears around the world every few years. Each time someone asks: what is the point of trying an old man who cannot even recognize his former colleagues? But the law has a reason worth stating fully. If "I have changed" were grounds for exoneration, every criminal could legally shed the past simply by waiting and changing enough; the justice victims await could never catch up with a changing person. The law makes a rough but deliberate choice: better a rough standard of identity than one so loose that anyone can slip free like a cicada leaving its shell. This reminds us that "the same person" is not merely a philosophical discovery. It is also a social convention, weighed down by interests and justice.

\textbf{The hospital room.} Medical decisions sharpen the question. When people with Alzheimer's, a deep coma or severe brain injury cannot express their wishes, who decides for them? The currently common practice is to respect "advance directives": a person writes while lucid, "If I reach such-and-such a condition, please treat me this way." Notice the philosophical judgment hidden here: \textbf{it assumes your lucid present self can legislate for a future self who may be completely different.} The institution partly accepts the uncle's reasoning: decisions should consult "the original person's" wishes, not merely the present responses of the body in the bed. Yet controversy persists. What if the future patient seems peaceful or even happy in that condition, while the lucid person once wrote, "Never let me live like that"? Who should be heard: the person who made the directive, or the person now breathing? These two "selves" sit opposite each other across time and illness, and the law must choose which to authorize.

\textbf{The phone.} The most everyday battleground is in your pocket. You have one or more online accounts containing conversations, photographs, game levels and networks of followers. These digital traces form an "informational continuity," like the continuity in Parfit's arrangement, and outlast your body. New questions line up: after someone dies, who owns the account? Is its content part of "him," or simply an inheritance, like a watch? Some platforms offer "memorial accounts," freezing the deceased person's timeline in place. This puts the relational view into technological practice: as long as a network of relationships remembers someone, they remain in some sense. A further step is "digital immortality" services, which train a conversational program on a person's entire chat history so a deceased loved one can keep "replying." Take this apart with the chapter's tools. It has informational continuity, perhaps mimicking very closely, but no bodily continuity. The relationship is one-way: you remember it; it is merely running. Is it him? Is it not him? Or is it a third kind of thing for which our language has no word? Your intuition stalls because our concept of identity was built for a world without these arrangements.

There is an even smaller, closer scene: your graduation autograph book. Almost every one says, "May you never change" or "Don't forget me." You now know how difficult those wishes are: never changing is impossible, and being forgotten is likely. Yet the writers are not foolish. What they are really writing may be a thread from the relational view: however we each change, \textbf{these shared days are already woven into our separate stories and cannot be removed}. Change is inevitable, but the past cannot be rewritten. Perhaps that is the only meaning of "forever" that can be fulfilled.

\section[Which "You" Would You Claim?]{For You: Which "You" Would You Be Willing to Claim?}
Return to the three things at the beginning: the child in the photograph, Grandma Chen in the hospital bed, and the "you" assembled from data on your phone.

This chapter has offered no baton that can never be dropped. The body's replacements leave only continuity; memories depart, personality changes direction, relationships let go, and stories admit several versions. Theseus's ship replaces every plank; the teletransporter's "you" walks the boundary between copying and destruction; the divided brain makes the word "identity" crash on the spot. Even awake, Zhuangzi cannot tell whether he is human or butterfly. Nagasena dismantles the chariot, and Parfit simply says there was no need to preserve the gold medal.

Now answer this question and give your reasons:

\textbf{If one day all your memories, ways of speaking and habits were copied completely into a machine, and its first words on waking were "I am you," would you acknowledge it as you?}

Before deciding, put all the weights on the scale. It has all your memories: the memory view should welcome it. It lacks your bodily continuity: the bodily view should reject it. It can maintain your relationships and care for those you love: the relational view and Parfit might ask what more you could want. But think again: after acknowledging it, does "you" still mean what it meant before? If it is you, what is the person reading this book now---the original, a copy, or one of two relics side by side?

Do not forget that the family in the ward still needs a course of action, not merely an account. Xiaoman will visit next Wednesday, and her grandmother will probably call her "Little Sister" again. Should Xiaoman correct her? Go along with her? Is following an older person's mistaken memories kindness, or complicity in abandoning the "real" grandmother? Is correcting her honest, or cruel?

There is no standard answer. But notice this: in answering "Is that me?" you use your memories, your body, your relationships and the story you tell yourself. What you use to vote is the ballot itself. Between the gap-toothed child in the photograph and you holding this book, nothing has stayed entirely unchanged. Yet this unbroken succession of changes is now thinking about itself. That fact comes closer to the heart of the mystery than any answer.
